% Root derivation: original EC three-column response, 19 September 2026.
\section{The joint response ellipse in the original three-column Gram}

This calculation starts from the same full retained eigenclass and the same
observation kernel as EC37--41. Its source is the supplied complete EC
continuation, preserved in \texttt{02\_Untitled 1775.md}, lines9401--10815;
the full-class completion is NCRP1--30. Every Gram below uses the original
arithmetic quotient metric $G_N$ and its full attained polynomial lift.
No native-section representative is substituted for the retained class.
The positive-square argument is the Hermitian Schur argument of
Boyd, El Ghaoui, Feron and Balakrishnan, \emph{Linear Matrix Inequalities
in System and Control Theory} (SIAM,1994), section2.1, equations(2.3)--(2.4),
doi:10.1137/1.9781611970777; its rank-deficient version is proved below
on the actual image, so no nonsingularity of a three-column Gram is assumed.

\subsection{Exact map, image, minimum and remainder}
Let $v=v_\lambda$ be the full retained class and let
$b_1=b_N$, $b_2=b_{N+1}$ be EC's original columns. Retain its original
kernel frame $I$, and write
\[
 P=I(I^*G_NI)^{-1}I^*G_N,\quad
 W=[b_1,b_2,v],\quad F=W^*G_NW,\quad K=W^*G_NPW .                 \tag{JRE1}
\]
The zero-dimensional kernel gives $P=0$ and the full kernel gives $P=1$.
In every case $P^2=P$ and $P^*G_N=G_NP$. In particular
$K=W^*P^*G_NPW$, so $0\preceq K\preceq F$.
Set
\[
 E=F_{33}>0,\quad f=\binom{F_{13}}{F_{23}},\quad
 \theta=K_{33}/E\in[0,1],\quad
 x=Pv,\quad z=x-\theta v .                                    \tag{JRE2}
\]
The two entries of $f$ are nonzero on EC's stated domain, by its original
Cauchy-transform proof EC20 and EC27--29. We keep their exact values.
Since $v^*G_Nx=x^*G_Nx=\theta E$, direct expansion gives
\[
 v^*G_Nz=0,\qquad z^*G_Nz=E\theta(1-\theta).                   \tag{JRE3}
\]
Let $B=[b_1,b_2]$ and define the exact residual-column map
\[
 A=B-vf^*/E:\mathbb C^2\longrightarrow E_{\rm packet},\quad
 S=A^*G_NA=F_{[1,2],[1,2]}-ff^*/E\succeq0 .                    \tag{JRE4}
\]
Here $E_{\rm packet}$ denotes the original polynomial quotient, while
$E$ in JRE2 is its retained-class squared norm. Thus their types differ
explicitly. The response residual is
\[
 r=A^*G_Nz=\binom{K_{13}}{K_{23}}-\theta f .                    \tag{JRE5}
\]
For every $c\in\ker S$ positivity implies $Ac=0$; hence $c^*r=0$.
Since $S$ is Hermitian, $r\in\operatorname{im}S$.
Define $S^\dagger$ as the inverse of $S$ on its positive eigenspaces,
and zero on its kernel. This specifies it at every rank without a
regularization parameter. The vector
$z_0=AS^\dagger r$ has $A^*G_Nz_0=r$, and $z-z_0$ is orthogonal to
$\operatorname{im}A$. Expanding the orthogonal square proves
\[
 \boxed{r^*S^\dagger r+
       \|z-AS^\dagger r\|_{G_N}^{2}=E\theta(1-\theta).}       \tag{JRE6}
\]
Thus $AS^\dagger r$ is the attained minimum-norm vector with the actual
two observations $A^*G_Nz=r$. Every other solution is its sum with a
vector in $\ker(A^*G_N)$, whose squared norm adds. This proves the
complete map, its range, its minimum, and its exact residual term.

To state the same identity in the original dimensionless responses,
retain EC37's $U=K_{13}/F_{13}$ and $V=K_{23}/F_{23}$ and put
\[
 D=\operatorname{diag}(F_{13},F_{23}),\quad
 C=D^{-1}A^*G_N,\quad T=D^{-1}S D^{-*},\quad
 w=Cz=\binom{U-\theta}{V-\theta}.                              \tag{JRE7}
\]
The codomain of $C$ is $\mathbb C^2$ with its standard Hermitian
metric; its adjoint from that space to the original $G_N$-space is
$C^\sharp=AD^{-*}$. In particular $CC^\sharp=T$.
The same orthogonal argument applied to $C$ proves
\[
 w\in\operatorname{im}T,\qquad
 w^*T^\dagger w+\|z-C^\sharp T^\dagger w\|_{G_N}^{2}
             =E\theta(1-\theta).                              \tag{JRE8}
\]
The two minimum vectors in JRE6 and JRE8 agree, since they solve
equivalent observations and both lie in $\operatorname{im}A$.
This proves the exact congruence of the image metrics even when
$S$ has rank zero or one. It does not use an invalid unrestricted
congruence formula for a Moore--Penrose inverse.

\subsection{A correlated finite bound for the three projected errors}
Write $t=T_{12}$, retain
\[
 L=E\theta(1-\theta),\quad
 d=\sqrt{L(T_{11}+T_{22}-2\Re t)},\quad
 h=\frac L2\left(|\Im t|+
                   \sqrt{T_{11}T_{22}-(\Re t)^2}\right).       \tag{JRE9}
\]
The quantities under square roots are nonnegative by $T\succeq0$.
All full source masses remain in $E$ and $T$; the products in JRE9
are independent of a common rescaling of the inner product only
because their displayed factors cancel exactly.
For $w=(u_1,u_2)^\mathsf T$, JRE8 and Cauchy--Schwarz give
\[
 |u_2-u_1|\le d,\qquad
 |u_j|^2\le L T_{jj}.                                        \tag{JRE10}
\]
Here is a sharper joint estimate for the quadratic term, including
its complex phase. Let
\[
 J=\begin{pmatrix}0&-i/2\\i/2&0\end{pmatrix}.
\]
Then $w^*Jw=\Im(\bar u_1u_2)$.
Write $w=T^{1/2}\eta$ with $\eta$ in the image of $T$;
JRE8 gives $\|\eta\|^2\le L$. The Hermitian matrix
$T^{1/2}JT^{1/2}$ has trace $-\Im t$ and determinant
$-\det(T)/4$. Its two eigenvalues are consequently
\[
 \frac{-\Im t\ \pm\sqrt{T_{11}T_{22}-(\Re t)^2}}2 .           \tag{JRE11}
\]
The formula remains true at rank zero or one, since trace and
determinant still determine its characteristic polynomial.
The spectral theorem now proves
\[
 |\Im(\bar u_1u_2)|\le h.                                   \tag{JRE12}
\]
This is the exact largest absolute quadratic value on the complete
ellipse JRE8 with its residual term omitted: an eigenvector of the
larger absolute eigenvalue attains it. The actual observation projection
need not attain that ellipse boundary; no such attainability is claimed.

Use EC38's three original products
\[
 P_K=\bar U V,\qquad P_B=(1-\bar U)(1-V),\qquad
 P_\times=\bar U+V-2\bar U V .                                \tag{JRE13}
\]
Substitute $U=\theta+u_1$, $V=\theta+u_2$. Their imaginary parts are
exactly
\[
 \begin{split}
 \Im P_K&=\theta\Im(u_2-u_1)+\Im(\bar u_1u_2),\\
 \Im P_B&=-(1-\theta)\Im(u_2-u_1)+\Im(\bar u_1u_2),\\
 \Im P_\times&=(1-2\theta)\Im(u_2-u_1)-2\Im(\bar u_1u_2).
 \end{split}                                                  \tag{JRE14}
\]
Their sum is zero. JRE10 and JRE12 prove the three simultaneous bounds
\[
 \begin{split}
 |\Im P_K|&\le\theta d+h,\\
 |\Im P_B|&\le(1-\theta)d+h,\\
 |\Im P_\times|&\le|1-2\theta|d+2h .
 \end{split}                                                  \tag{JRE15}
\]
All three use the same $T$, the same $\theta$, and the same retained
class. In particular a small eigenvalue or a small original $F_{j3}$
is visible through the exact entries of $T$; it is never removed
by replacing the Gram with an identity or a rank fraction.

In EC39 let $\varepsilon_{N,\lambda}>0$ be its single shared phase
remainder and $E_{\rm cur}$ its proved finite bound. The exact
response-error vector, in the order $(K,B,\times)$, is
\[
 \mathcal R_X=2E(-1)^r\varepsilon_{N,\lambda}\Im P_X,
 \qquad \sum_X\mathcal R_X=0 .                               \tag{JRE16}
\]
Thus its individual absolute values are bounded, respectively, by
\[
 2EE_{\rm cur}(\theta d+h),\quad
 2EE_{\rm cur}((1-\theta)d+h),\quad
 2EE_{\rm cur}(|1-2\theta|d+2h).                              \tag{JRE17}
\]
This is a finite amplification estimate on the actual attained
subquotient metrics. It leaves the actual two-column Schur complement
and the original period-dependent kernel projection to be evaluated.
It does not assign a sign to an individual projected term from the
positive full-class sum. The next analytic task is now the joint
profile of $T$ and $\theta$ in these exact original coordinates.

\subsection{The actual parity map and its two boundary energies}
For the original even arithmetic measure on $S=a/2+iy$, reflection
$\mathcal R P(S)=P(a-S)$ is an isometry. The full root polynomial
satisfies $\chi(a-S)=(-1)^q\chi(S)$, so reflection preserves each
complete relation space $\chi\mathbb C[S]_{\le N-q}$. It therefore
induces an isometry of the attained quotient, commuting with its
minimum lift. The two original boundary columns have opposite
parities $(-1)^N$ and $(-1)^{N+1}$: their source monic orthogonal
polynomials do, by uniqueness of the monic minimum in the even
measure, and the remainder map commutes with reflection. Consequently
\[
 F_{12}=\langle b_N,b_{N+1}\rangle_{G_N}
       =-\langle b_N,b_{N+1}\rangle_{G_N}=0 .                   \tag{JRE18}
\]
This is the full quotient identity also calculated in NCE18--21.
It preserves the actual source, period, conductor and class. Inserting
it in JRE4 and JRE7 gives the exact original response matrix
\[
 T=\begin{pmatrix}F_{11}/|F_{13}|^2&0\\0&F_{22}/|F_{23}|^2\end{pmatrix}
       -\frac1E\begin{pmatrix}1&1\\1&1\end{pmatrix},\qquad
 T_{12}=-1/E .                                                \tag{JRE19}
\]
Define the two actual boundary energies
\[
 p_1=\frac{|F_{13}|^2}{EF_{11}},\qquad
 p_2=\frac{|F_{23}|^2}{EF_{22}}.                              \tag{JRE20}
\]
Both are positive. The two columns $b_j/\sqrt{F_{jj}}$ are
orthonormal by JRE18. The orthogonal square of their projection
of $v/\sqrt E$ is $p_1+p_2$, so $p_1+p_2\le1$ with the exact
remainder equal to the squared norm of its orthogonal complement.
In particular these are values of specified original projections,
not fractions assigned from ranks.
Substitute JRE19--20 into JRE9. Every original factor cancels
only inside its displayed matched ratio, yielding
\[
 \boxed{d^2=\theta(1-\theta)(p_1^{-1}+p_2^{-1}),\qquad
 h=\frac{\theta(1-\theta)}{2\sqrt{p_1p_2}}
                    \sqrt{1-p_1-p_2}.}                       \tag{JRE21}
\]
Together JRE17 and JRE21 reduce the actual amplification estimate
to the three original scalar projection energies
$p_1,p_2,\theta$. In particular the exact total allowance obeys
\[
 \sum_{X=K,B,\times}|\mathcal R_X|
 \le2EE_{\rm cur}\{(1+|1-2\theta|)d+4h\},\qquad
 \sum_X\mathcal R_X=0.                                      \tag{JRE22}
\]
The first statement follows by adding the three proved absolute
bounds, while the second is the unchanged signed identity JRE16.
No lower bound for either $p_j$ is inserted. Their possible smallness
remains explicitly in JRE21 and is a concrete target for the next
original-source calculation.
